\documentclass[11pt]{article}
\usepackage[margin=0.9in]{geometry}
\usepackage{amsmath,amssymb,amsthm,mathtools}
\usepackage{graphicx}
\usepackage{booktabs}
\usepackage[hidelinks]{hyperref}
\usepackage{microtype}
\setlength{\emergencystretch}{2em}

\newtheorem{theorem}{Theorem}
\newtheorem{lemma}[theorem]{Lemma}
\newtheorem{proposition}[theorem]{Proposition}
\newcommand{\K}{\mathbb K}
\newcommand{\N}{\mathbb N}
\newcommand{\Lcal}{\mathcal L}
\newcommand{\ellw}{\ell^{w}}

\title{A Complete Low-Exponent Anisotropic\\
Grothendieck Coincidence Theorem\\
\large and a High-Exponent Hadamard Obstruction}
\author{Automated mathematics run \texttt{fa\_banach\_001}, lane 2}
\date{26 August 2026}

\begin{document}
\maketitle

\begin{center}
\fbox{\parbox{0.91\textwidth}{
\textbf{Status: substantial partial result likely valid.}
For arbitrary anisotropic input exponents in \([1,2]\), we give a complete
if-and-only-if classification.  The source's unrestricted high-exponent
classification remains open, but a new ordered pairwise obstruction rules
out the naive coordinatewise extension.}}
\end{center}

\begin{abstract}
Albuquerque, Ara\'ujo, Rezende, and Santos ask for which anisotropic
multi-indices every bounded \(m\)-linear map
\(\ell_1^m\to\ell_2\) is multiple summing.  We completely answer the
question when all input exponents lie in \([1,2]\): for
\(\mathbf p=(p_1,\ldots,p_m)\in[1,2]^m\) and
\(\mathbf q=(q_1,\ldots,q_m)\in[1,\infty)^m\), the coincidence holds if
and only if \(q_k\ge p_k\) for every \(k\).  At the vertices
\(p_k\in\{1,2\}\), mixed-norm Minkowski reorders the exponent-one levels
outside the exponent-two levels.  The inner block is controlled by the
multilinear \((2;2)\) Grothendieck theorem and the outer block by the
\((1;1)\) theorem, both with Hilbert-valued range.  Coordinatewise complex
interpolation fills the cube \([1,2]^m\), and mixed-norm monotonicity gives
all \(q_k\ge p_k\).  Rank-one maps prove necessity.  Beyond this cube, a
Hadamard construction shows that every coincidence tuple must satisfy
\[
 \frac1{q_k}\le \frac1{p_\ell}
 +\left(\frac1{p_k}-\frac12\right)_+
 \qquad (k<\ell).
\]
Thus the unrestricted problem has genuinely ordered cross-coordinate
constraints; coordinatewise scalar conditions cannot classify it.
\end{abstract}

\section{Source question}

The source is N.~Albuquerque, G.~Ara\'ujo, L.~Rezende, and J.~Santos,
\emph{A summability principle and applications}, arXiv:1904.04549v2
(2025)\cite{source}.  On PDF page 9 it asks for which anisotropic
multi-indices \((\mathbf q;\mathbf p)\) one has
\[
 \Pi^{\mathrm{ms}}_{(\mathbf q;\mathbf p)}
   ({}^m\ell_1;\ell_2)=\Lcal({}^m\ell_1;\ell_2).
\]
The source recalls the complete isotropic answer and supplies additional
sufficient anisotropic families, but not a full classification.

\begin{figure}[ht]
  \centering
  \includegraphics[width=0.94\textwidth]{figures/open_problem_crop.png}
  \caption{The anisotropic classification question and the source's
  Proposition~4.7, PDF page 9.}
\end{figure}

For Banach-space-valued arrays \(a=(a_{i_1,\ldots,i_m})\), write
\begin{align*}
 \|a\|_{\ell_{\mathbf q}}
 :=\left(\sum_{i_1}
 \left(\cdots\left(\sum_{i_m}\|a_{i_1,\ldots,i_m}\|^{q_m}
 \right)^{q_{m-1}/q_m}\cdots\right)^{q_1/q_2}
 \right)^{1/q_1}.
\end{align*}
An \(m\)-linear \(T:\ell_1^m\to\ell_2\) is multiple
\((\mathbf q;\mathbf p)\)-summing when
\begin{equation}\label{eq:def}
 \big\|(T(x_{i_1}^1,\ldots,x_{i_m}^m))_{\mathbf i}\big\|_{\ell_{\mathbf q}}
 \le C\prod_{k=1}^m\|(x_i^k)_i\|_{w,p_k}
\end{equation}
for every finitely supported choice of sequences.

\section{Proof intuition}

The isotropic theorem gives the two endpoints needed for each coordinate:
every \(m\)-linear map \(\ell_1^m\to H\), with \(H\) Hilbert, is both
multiple \((1;1)\)- and multiple \((2;2)\)-summing.  A mixed vertex such
as \((1,2,1,2)\) appears different, but Minkowski's inequality moves the
smaller exponents outward, where the mixed norm is largest.  After this
reordering, the exponent-two coordinates form one Hilbert-valued inner
block.  The exponent-one coordinates then act on that Hilbert space as an
outer block.  Applying the endpoint theorem twice proves every
\(\{1,2\}\)-vertex.

The remaining step is genuinely anisotropic: interpolate one coordinate at
a time, simultaneously interpolating its weak sequence space and the
corresponding level of the mixed output norm.  The required uniform identity
for weak sequence spaces over \(\ell_1\) is the standard injective-tensor
interpolation formula used in Bayart--Pellegrino--Rueda\cite{bpr}.  This
fills the entire low-exponent cube.

\section{The exact low-exponent theorem}

\begin{theorem}\label{thm:main}
Let \(m\ge1\), let \(H\) be a Hilbert space, and let
\[
 \mathbf p=(p_1,\ldots,p_m)\in[1,2]^m,
 \qquad \mathbf q=(q_1,\ldots,q_m)\in[1,\infty)^m.
\]
Then
\[
 \Pi^{\mathrm{ms}}_{(\mathbf q;\mathbf p)}({}^m\ell_1;H)
 =\Lcal({}^m\ell_1;H)
\]
if and only if \(q_k\ge p_k\) for every \(k\).
In particular, this completely classifies the source problem throughout
the input cube \([1,2]^m\).
\end{theorem}

\begin{proof}
We first prove sufficiency when \(q_k=p_k\in\{1,2\}\) for every \(k\).
For a nonnegative matrix and \(1\le r\le s\), Minkowski's inequality gives
\begin{equation}\label{eq:minkowski}
 \left(\sum_i\left(\sum_j a_{ij}^{r}\right)^{s/r}\right)^{1/s}
 \le
 \left(\sum_j\left(\sum_i a_{ij}^{s}\right)^{r/s}\right)^{1/r}.
\end{equation}
Thus each adjacent pattern with exponent \(2\) outside exponent \(1\) is
bounded by the transposed pattern with exponent \(1\) outside exponent
\(2\).  Repeatedly applying \eqref{eq:minkowski}, while permuting the
corresponding variables of \(T\), reduces every vertex to
\[
 (\underbrace{1,\ldots,1}_{a\text{ times}},
  \underbrace{2,\ldots,2}_{b\text{ times}}),\qquad a+b=m.
\]

Fix sequences in the last \(b\) variables.  Define
\[
 S(x_1,\ldots,x_a)
 =\big(T(x_1,\ldots,x_a,x^{a+1}_{i_{a+1}},\ldots,
 x^m_{i_m})\big)_{i_{a+1},\ldots,i_m}.
\]
The target \(\ell_2(H)\) over the last \(b\) indices is Hilbert.  The
multilinear \((2;2)\) Grothendieck coincidence, which is the \(p=2\)
case of source Theorem~4.6, gives
\[
 \|S\|\le C_{2,b}\|T\|
 \prod_{k=a+1}^m\|(x_i^k)_i\|_{w,2}.
\]
Now apply the \((1;1)\) coincidence to the \(a\)-linear map \(S\), again
with Hilbert target.  We obtain
\begin{align*}
 &\sum_{i_1,\ldots,i_a}
 \left(\sum_{i_{a+1},\ldots,i_m}
 \|T(x^1_{i_1},\ldots,x^m_{i_m})\|^2\right)^{1/2}\\
 &\qquad\le C_{1,a}C_{2,b}\|T\|
 \prod_{k=1}^{a}\|(x_i^k)_i\|_{w,1}
 \prod_{k=a+1}^{m}\|(x_i^k)_i\|_{w,2}.
\end{align*}
This is exactly the sorted mixed-vertex estimate.  The cases \(a=0\) or
\(b=0\) use only the corresponding endpoint theorem.

We next fill the cube.  For a finite sequence of length \(N\), put
\[
 W_p^N=\ell_p^{w,N}(\ell_1).
\]
The canonical identification
\[
 W_p^N=\ell_p^N\widehat\otimes_\varepsilon\ell_1
\]
is isometric.  Uniformly in \(N\), injective-tensor interpolation gives an
isomorphism
\begin{equation}\label{eq:weak-interp}
 [W_1^N,W_2^N]_\theta\simeq W_p^N,
 \qquad \frac1p=1-\frac\theta2.
\end{equation}
This is the \(\ell_1\) weak-sequence interpolation formula used in
\cite[Section~2]{bpr}; it follows from the injective interpolation theorem
there.  Vector-valued sequence-space interpolation also gives, at any fixed
level of an iterated mixed norm,
\[
 [\ell_1(Y),\ell_2(Y)]_\theta=\ell_p(Y).
\]

For fixed \(T\) and \(N\), its coefficient map is multilinear from the
product of the \(W_{p_k}^N\) spaces to the corresponding finite mixed
sequence space.  We have just proved a uniform bound at all \(2^m\) cube
vertices.  Apply multilinear complex interpolation first in coordinate 1,
then coordinate 2, and so on, using \eqref{eq:weak-interp} and the
vector-valued output formula at each step.  This proves \eqref{eq:def} with
\(\mathbf q=\mathbf p\) for every \(\mathbf p\in[1,2]^m\), with a constant
independent of \(N\).  Taking the supremum over finite truncations gives the
infinite-sequence assertion.  For real scalars, complexification gives the
same coincidence with a harmless constant depending only on \(m\).

If \(q_k\ge p_k\) for every \(k\), coordinatewise monotonicity of iterated
\(\ell_p\) norms yields
\(\|a\|_{\ell_{\mathbf q}}\le\|a\|_{\ell_{\mathbf p}}\), proving
sufficiency in full.

For necessity, fix \(k\).  Choose norm-one functionals
\(\varphi_j\in\ell_1^*=\ell_\infty\), a unit vector \(h\in H\), and the
rank-one map
\[
 T(x_1,\ldots,x_m)=\varphi_1(x_1)\cdots\varphi_m(x_m)h.
\]
In every coordinate except \(k\), use a one-term sequence on which
\(\varphi_j\) equals one.  In coordinate \(k\), use
\(x_i=a_i u\) with \(\varphi_k(u)=1\).  Inequality \eqref{eq:def} reduces to
\[
 \|(a_i)\|_{q_k}\le C\|(a_i)\|_{p_k}.
\]
Uniformity over finite scalar sequences forces \(q_k\ge p_k\).  This holds
for every \(k\), completing the proof.
\end{proof}

\section{A high-exponent Hadamard obstruction}

The low-exponent theorem does not extend coordinatewise once a later input
exponent is allowed above two.  The following necessary condition makes the
new obstruction explicit.

\begin{proposition}\label{prop:hadamard}
Assume \(m\ge2\) and
\[
 \Pi^{\mathrm{ms}}_{(\mathbf q;\mathbf p)}({}^m\ell_1;H)
 =\Lcal({}^m\ell_1;H)
\]
for a nonzero Hilbert space \(H\).  Then, besides \(q_j\ge p_j\) for every
\(j\), one necessarily has
\begin{equation}\label{eq:hadamard-obstruction}
 \frac1{q_k}\le \frac1{p_\ell}
 +\max\left\{\frac1{p_k}-\frac12,0\right\}
 \qquad\text{for every }1\le k<\ell\le m.
\end{equation}
\end{proposition}

\begin{proof}
Equality of the two operator spaces and the open mapping theorem give a
constant \(C\), independent of \(T\), in \eqref{eq:def}.  Fix \(k<\ell\)
and let \(N\) run through powers of two.  Let
\(A_N=(\varepsilon_{ab})_{a,b=1}^N\) be a Sylvester Hadamard matrix, so
\(A_NA_N^{\mathsf T}=NI_N\).  On the \(k\)-th copy of \(\ell_1\), put
\[
 x_i^k=\frac1N\sum_{a=1}^N\varepsilon_{ia}e_a,
 \qquad i=1,\ldots,N,
\]
and on the \(\ell\)-th copy put \(x_j^\ell=e_j\).  For every other
coordinate use a single unit vector.

If \(z\in B_{\ell_\infty^N}\), orthogonality gives
\(\|A_Nz\|_2=\sqrt N\|z\|_2\le N\).  Comparing finite-dimensional
\(\ell_p\)-norms therefore yields
\begin{equation}\label{eq:hadamard-weak}
 \big\|(x_i^k)_{i=1}^N\big\|_{w,p_k}
 =\frac1N\|A_N:\ell_\infty^N\to\ell_{p_k}^N\|
 \le N^{(1/p_k-1/2)_+}.
\end{equation}
Also \(\|(e_j)_{j=1}^N\|_{w,p_\ell}=N^{1/p_\ell}\).

Choose a unit vector \(h\in H\), norm-one functionals in the unused
variables, and define \(T_N\) by inserting in coordinates \(k,\ell\) the
bilinear form
\[
 B_N(u,v)=\sum_{a,b=1}^N\varepsilon_{ab}u_av_b
\]
and multiplying by \(h\) and those functionals.  Since the domain is a
product of \(\ell_1\)-spaces, \(\|T_N\|=\max_{a,b}|\varepsilon_{ab}|=1\).
Hadamard orthogonality gives
\[
 T_N(x_i^k,e_j)=\delta_{ij}h.
\]
All intervening indices have length one, so the mixed norm of this diagonal
array is \(N^{1/q_k}\).  Applying \eqref{eq:def} and
\eqref{eq:hadamard-weak} gives
\[
 N^{1/q_k}\le C
 N^{(1/p_k-1/2)_++1/p_\ell}.
\]
Letting \(N\to\infty\) proves \eqref{eq:hadamard-obstruction}.
\end{proof}

For example, \(\mathbf p=(3,100)\) and \(\mathbf q=(4,101)\) satisfy the
naive coordinatewise inequalities, including strict improvement in both
high coordinates, but Proposition~\ref{prop:hadamard} would require
\(1/4\le1/100\).  Hence this tuple is not a coincidence tuple.

\section{Verification, novelty, and remaining range}

The script \texttt{code/verify\_mixed\_norm\_reductions.py} checks the
adjacent Minkowski direction, bubble-sort reduction of arbitrary
\(\{1,2\}\)-patterns, and coordinatewise mixed-norm monotonicity on random
arrays in dimensions two through five.  It also checks the Sylvester
orthogonality and weak-norm estimate used in Proposition~\ref{prop:hadamard}.
These computations guard against norm-order and scaling mistakes; they are
not part of the proof.

The theorem is sharp and complete for \(p_k\le2\).  The unrestricted source
question remains open here.  Freezing all other variables gives only the
coordinatewise necessity \(q_k\ge p_k\).  Proposition~\ref{prop:hadamard}
shows that this is far from sufficient: later input exponents constrain
earlier output levels.  A complete sufficiency theorem for these ordered
high-exponent interactions is not supplied by the known isotropic theorem or
the source's suffix-balance inclusion families.

A bounded novelty audit on 26 August 2026 searched all four run indexes, the
local parsed arXiv corpus, and arXiv/web results using the source id/title and
the phrases ``anisotropic Grothendieck,'' ``mixed norm multiple summing,''
and ``\(\ell_1\) to \(\ell_2\) anisotropic coincidence.''  It found the
source, the earlier isotropic and interpolation papers, and no statement of
the complete \([1,2]^m\) classification above.  The source's 2025 revision
still poses the general question.  This is bounded evidence, not a priority
claim.

Human review recommendation: \textbf{send to human}, focusing on the
coordinatewise interpolation step, uniformity of \eqref{eq:weak-interp}, and
the exponent direction in \eqref{eq:hadamard-obstruction}.

\enlargethispage{2\baselineskip}
\begin{thebibliography}{9}
\footnotesize
\raggedright
\bibitem{source}
N.~Albuquerque, G.~Ara\'ujo, L.~Rezende, and J.~Santos,
\emph{A summability principle and applications}, arXiv:1904.04549v2, 2025.

\bibitem{bpr}
F.~Bayart, D.~Pellegrino, and P.~Rueda,
\emph{On coincidence results for summing multilinear operators:
interpolation, \(\ell_1\)-spaces and cotype}, arXiv:1805.12500, 2018.
\end{thebibliography}

\end{document}
